% Historia

%%%%%%%%%%%%%%%%%%%%%%%%%%%%%%%%%%%%%%%%%%%%%%%%%%%%%%%%%%%%%%%%%%%%%%%%%%%%%
%%%%%%%%%%%%%%%%%%%%%%%%%%%%%%%%%%%%%%%%%%%%%%%%%%%%%%%%%%%%%%%%%%%%%%%%%%%%%
%%%%%%%%%%%%%% ASISTENCIA PARA REDACCIÓN %%%%%%%%%%%%%%%%%%%%%%%%%%%%%%%%%%%%
%%%%%%%%%%%%%%%%%%%%%%%%%%%%%%%%%%%%%%%%%%%%%%%%%%%%%%%%%%%%%%%%%%%%%%%%%%%%%

% \cultivo -> genera un lienzo para 1 dibujo de tamaño normal
% \cultivos -> genera un lienzo para 2 o más dibujos de tamaño normal
% \cultivoP -> genera un lienzo para 1 dibujo de tamaño pequeño
% \cultivosP -> genera un lienzo para 2 o más dibujos de tamaño pequeño
% \granja -> dibuja una granja hexagonal adecuada al tamaño del lienzo
% \granjaN -> dibuja una granja hexagonal adecuada al tamaño del lienzo; el hexágono está relleno de color negro
% \fin -> finaliza un lienzo de 1 dibujo
% \finfin -> finaliza un lienzo de 2 o más dibujos

% node distance en tamaño normal = 10.2mm
% node distance en tamaño pequeño = 7.8mm
% arriba derecha = 30
% abajo derecha = 330
% arriba izquierda = 150
% abajo izquierda = 210
% arriba = [above of = #]      ; # es el nombre del nodo de referencia
% abajo = [below of = #]       ; # es el nombre del nodo de referencia

%%%%%%%%%%%%%%%%%%%%%%%%%%%%%%%%%%%%%%%%%%%%%%%%%%%%%%%%%%%%%%%%%%%%%%%
%%%%%%%%%%%%%%%%%%%%%%%%%%%%%  EJEMPLOS %%%%%%%%%%%%%%%%%%%%%%%%%%%%%%%
%%%%%%%%%%%%%%%%%%%%%%%%%%%%%%%%%%%%%%%%%%%%%%%%%%%%%%%%%%%%%%%%%%%%%%%

%\cultivo
%\granja (g1) {$G_1$} ;
%\granja (g2) at ($ (g1) + (330:10.2mm) $) {$G_2$} ; % abajo derecha
%\granja (g3) at ($ (g1) + (210:10.2mm) $) {$G_3$} ; % abajo izquierda
%\granja (g4) at ($ (g1) + (150:10.2mm) $) {$G_4$} ; % arriba izquierda
%\granja (g5) at ($ (g1) + (30:10.2mm) $) {$G_5$} ; % arriba derecha
%\granja (g6) [above of = g1] {$G_6$} ; % arriba
%\granja (g7) [below of = g1] {$G_7$} ; % abajo
%\fin

%\vspace{10.2mm}

%\cultivoP
%\granja (g1) {$G_1$} ;
%\granja (g2) at ($ (g1) + (330:7.8mm) $) {$G_2$} ; % abajo derecha
%\granja (g3) at ($ (g1) + (210:7.8mm) $) {$G_3$} ; % abajo izquierda
%\granja (g4) at ($ (g1) + (150:7.8mm) $) {$G_4$} ; % arriba izquierda
%\granja (g5) at ($ (g1) + (30:7.8mm) $) {$G_5$} ; % arriba derecha
%\granja (g6) [above of = g1] {$G_6$} ; % arriba
%\granja (g7) [below of = g1] {$G_7$} ; % abajo
%\fin

%\vspace{10.2mm}

%\begin{center}
%\cultivos
%\granja (g1) {$G_1$} ;
%\granja (g2) at ($ (g1) + (330:10.2mm) $) {$G_2$} ; % abajo derecha
%\granja (g3) at ($ (g1) + (210:10.2mm) $) {$G_3$} ; % abajo izquierda
%\granja (g4) at ($ (g1) + (150:10.2mm) $) {$G_4$} ; % arriba izquierda
%\granja (g5) at ($ (g1) + (30:10.2mm) $) {$G_5$} ; % arriba derecha
%\granja (g6) [above of = g1] {$G_6$} ; % arriba
%\granja (g7) [below of = g1] {$G_7$} ; % abajo
%\finfin
%\hspace{2mm}
%\cultivos
%\granja (g1) {$G_1$} ;
%\granja (g2) at ($ (g1) + (330:10.2mm) $) {$G_2$} ; % abajo derecha
%\granja (g3) at ($ (g1) + (210:10.2mm) $) {$G_3$} ; % abajo izquierda
%\granja (g4) at ($ (g1) + (150:10.2mm) $) {$G_4$} ; % arriba izquierda
%\granja (g5) at ($ (g1) + (30:10.2mm) $) {$G_5$} ; % arriba derecha
%\granja (g6) [above of = g1] {$G_6$} ; % arriba
%\granja (g7) [below of = g1] {$G_7$} ; % abajo
%\finfin
%\end{center}

%\vspace{10.2mm}

%\begin{center}
%\cultivosP
%\granja (g1) {$G_1$} ;
%\granja (g2) at ($ (g1) + (330:7.8mm) $) {$G_2$} ; % abajo derecha
%\granja (g3) at ($ (g1) + (210:7.8mm) $) {$G_3$} ; % abajo izquierda
%\granja (g4) at ($ (g1) + (150:7.8mm) $) {$G_4$} ; % arriba izquierda
%\granja (g5) at ($ (g1) + (30:7.8mm) $) {$G_5$} ; % arriba derecha
%\granja (g6) [above of = g1] {$G_6$} ; % arriba
%\granja (g7) [below of = g1] {$G_7$} ; % abajo
%\finfin
%\hspace{2mm}
%\cultivosP
%\granja (g1) {$G_1$} ;
%\granja (g2) at ($ (g1) + (330:7.8mm) $) {$G_2$} ; % abajo derecha
%\granja (g3) at ($ (g1) + (210:7.8mm) $) {$G_3$} ; % abajo izquierda
%\granja (g4) at ($ (g1) + (150:7.8mm) $) {$G_4$} ; % arriba izquierda
%\granja (g5) at ($ (g1) + (30:7.8mm) $) {$G_5$} ; % arriba derecha
%\granja (g6) [above of = g1] {$G_6$} ; % arriba
%\granja (g7) [below of = g1] {$G_7$} ; % abajo
%\finfin
%\end{center}

%%%%%%%%%%%%%%%%%%%%%%%%%%%%%%%%%%%%%%%%%%%%%%%%%%%%%%%%%%%%%%%%%%%%%%%%%%%
%%%%%%%%%%%%%%%%%%%%% PRUEBA %%%%%%%%%%%%%%%%%%%%%%%%%%%%%%%%%%%%%%%%%%%%%%
%%%%%%%%%%%%%%%%%%%%%%%%%%%%%%%%%%%%%%%%%%%%%%%%%%%%%%%%%%%%%%%%%%%%%%%%%%%

% historia
\noindent Tras unas breves vacaciones en las lunas de Júpiter, el equipo académico de CMAT ha %
sido designado para emprender una nueva misión. Liderados por Patito, deberán explorar la galaxia %
en busca de un planeta con condiciones para el desarrollo de vida humana. De esta forma, Patito y %
su equipo de aprendices se han embarcado en la nave \emph{Conejo Matemático Misterioso 1} o \emph{CMM-1} %
y han comenzado un largo y emocionante viaje hacia los confines del universo.

\noindent Una tranquila jornada en el CMM-1 se vivía rodeada de café y galletas matemáticas. Patito, %
sentado en una de las sillas del casino de la nave, disfrutaba de un delicioso completo cuando un grito interrumpió %
su descanso para comer. --- \emph{¡Lo encontramos!} --- se escuchó repetidamente mientras las pantallas de la nave %
mostraban los resultados del análisis de un planeta. El planeta en cuestión, era habitable. Patito, al borde de las %
lágrimas de emoción, decidió nombrar al planeta como \emph{Champiñón}.

\noindent Después de comunicar las excelentes noticias a la Tierra, Patito recibió una orden incluso más desafiante %
que la anterior: debía aterrizar en el planeta, explorarlo, hacer mapas y prepararlo para la llegada de más humanos, %
con el propósito de crear una colonia en él. Patito llamaría a esta colonia como; \emph{Colonia Champiñón Primera}.

\section*{Granjas Espaciales}

\noindent Uno de los desafíos más importantes de Patito, sería el de generar comida para la llegada de nuevos humanos al %
planeta. Como este era habitable, poseía una atmósfera similar a la Tierra y un suelo que se podría considerar fértil, %
Patito decidió presentar al equipo académico su nuevo invento.

\noindent Patito llamaría a este invento como \textbf{\emph{Granjas Espaciales}}, a las que designaremos por $\mathbf{G}$. %
Estas granjas espaciales son estructuras hexagonales que pueden generar comida y se pueden unir entre ellas. Cuando %
\textbf{dos o más granjas espaciales} se unen entre ellas, a este grupo le llamamos \textbf{\emph{Campo de Cultivo}}. En los %
diagramas inferiores se pueden ver \emph{campos de cultivo} con $2$, $3$ y $4$ \emph{granjas espaciales} respectivamente.

\begin{center}
\cultivos
\granja (g1) {$G_1$} ;
\granja (g2) at ($(g1) + (30:\grande)$) {$G_2$} ;
\finfin
\hspace{20mm}
\cultivos
\granja (g1) {$G_1$} ;
\granja (g2) at ($(g1) + (30:\grande)$) {$G_2$} ;
\granja (g3) at ($(g1) + (330:\grande)$) {$G_3$} ;
\finfin
\hspace{20mm}
\cultivos
\granja (g1) {$G_1$} ;
\granja (g2) at ($(g1) + (30:\grande)$) {$G_2$} ;
\granja (g3) at ($(g1) + (330:\grande)$) {$G_3$} ;
\granja (g4) at ($(g3) + (30:\grande)$) {$G_4$} ;
\finfin
\end{center}

\noindent Mientras Patito señalaba al equipo académico la importancia de estas granjas espaciales en la generación de comida %
para las futuras colonias, explicó los detalles cruciales en el desarrollo de su idea. Las granjas cambiarían su forma de %
generar comida, dependiendo de la forma en que estas estén unidas entre si. Mencionaremos por ahora dos tipos de uniones.

\begin{itemize}
\item Conoceremos como \textbf{\emph{granjas vecinas}} a todas aquellas granjas que \textbf{comparten una arista} entre si. Una %
definición alternativa corresponde a: Son \emph{granjas vecinas} de una granja espacial, todas aquellas granjas que se encuentren %
adyacentes a esta.
\item Conoceremos como \textbf{\emph{supergranja}} al conjunto de exactamente tres granjas espaciales que \textbf{compartan un vértice}.
\end{itemize}

\vspace{2mm}

\noindent En los ejemplos inferiores podemos evaluar la presencia de \emph{granjas vecinas}. En el diagrama de la izquierda, $G_1$ tiene %
una granja vecina ($G_2$) y $G_2$ tiene una granja vecina ($G_1$). En el diagrama de la derecha, $G_1$, $G_2$ y $G_3$ tienen todas dos granjas %
vecinas ($G_2$ y $G_3$), ($G_1$ y $G_3$) y ($G_1$ y $G_2$) respectivamente. Las \emph{supergranjas}, por otro lado, están únicamente presentes %
en el diagrama de la derecha. En este diagrama encontramos un total de una supergranja, comformada por las granjas espaciales $G_1$, $G_2$ y $G_3$.

\begin{center}
\cultivos
\granja (g1) {$G_1$} ;
\granja (g2) at ($(g1) + (30:\grande)$) {$G_2$} ;
\finfin
\hspace{20mm}
\cultivos
\granja (g1) {$G_1$} ;
\granja (g2) at ($(g1) + (30:\grande)$) {$G_2$} ;
\granja (g3) at ($(g1) + (330:\grande)$) {$G_3$} ;
\finfin
\end{center}

\vspace{2mm}

\noindent Notemos también que una granja espacial puede ser parte de más de una \emph{supergranja}. En el siguiente diagrama existen dos supergranjas, %
conformadas por las granjas ($G_1$, $G_2$ y $G_3$) y ($G_2$, $G_3$ y $G_4$) respectivamente.


\cultivo
\granja (g1) {$G_1$} ;
\granja (g2) at ($(g1) + (30:\grande)$) {$G_2$} ;
\granja (g3) at ($(g1) + (330:\grande)$) {$G_3$} ;
\granja (g4) at ($(g3) + (30:\grande)$) {$G_4$} ;
\fin

\subsection*{Generación de alimentos; Reglas}

Una vez el equipo académico entendió cómo funcionaba el invento de Patito, llegó el momento de resolver una pregunta crucial. ¿Cuánto alimento %
genera una granja espacial? Para responder ello, Patito presentó al equipo un set de tres reglas que definen la generación de alimentos de una granja espacial %
en cada caso. Es importante señalar que al momento de calcular la cantidad de alimento que genera una granja, cada una de estas reglas deberá ser verificada %
de manera \textbf{independiente}.

\begin{itemize}
\item Cada \emph{granja espacial} produce una cantidad base de $\mathbf{2}$ \textbf{Alimentos}. \textbf{(REGLA 1)}
\item Cada granja espacial obtiene un adicional de $\mathbf{1}$ \textbf{Alimento} por \textbf{cada} \emph{granja vecina} que esta posea. \textbf{(REGLA 2)}
\item Cada granja espacial obtiene un adicional de $\mathbf{2}$ \textbf{Alimentos} por \textbf{cada} \emph{supergranja} a la que esta pertenezca. \textbf{(REGLA 3)}
\end{itemize}

\noindent Como ejemplo evaluaremos la cantidad de alimento generada por el siguiente campo de cultivo, compuesto por las granjas $G_1$, $G_2$ y $G_3$. Para ello, %
evaluaremos a cada granja de manera independiente.

\cultivo
\granja (g1) {$G_1$} ;
\granja (g2) at ($(g1) + (30:\grande)$) {$G_2$} ;
\granja (g3) at ($(g1) + (330:\grande)$) {$G_3$} ;
\fin

\begin{itemize}
\item La granja espacial $G_1$ obtiene $2$ unidades de alimento por \emph{Regla 1}. Sus granjas vecinas son $G_2$ y $G_3$, en total, $2$, por lo que obtendrá un adicional %
de $2$ de alimento por \emph{Regla 2}. También podemos notar que pertenece a la supergranja compuesta por $G_1$, $G_2$ y $G_3$, obteniendo así un adicional de $2$ de alimento %
por \emph{Regla 3}.
\item La granja espacial $G_2$ obtiene $2$ unidades de alimento por \emph{Regla 1}. Sus granjas vecinas son $G_1$ y $G_3$, en total, $2$, por lo que obtendrá un adicional %
de $2$ de alimento por \emph{Regla 2}. También podemos notar que pertenece a la supergranja compuesta por $G_1$, $G_2$ y $G_3$, obteniendo así un adicional de $2$ de alimento %
por \emph{Regla 3}.
\item La granja espacial $G_3$ obtiene $2$ unidades de alimento por \emph{Regla 1}. Sus granjas vecinas son $G_2$ y $G_1$, en total, $2$, por lo que obtendrá un adicional %
de $2$ de alimento por \emph{Regla 2}. También podemos notar que pertenece a la supergranja compuesta por $G_1$, $G_2$ y $G_3$, obteniendo así un adicional de $2$ de alimento %
por \emph{Regla 3}.
\end{itemize}

\noindent Otra forma de desarrollar este caso podría haber sido notar que las tres granjas son equivalentes u homólogas, generando todas la misma cantidad de alimento. Por ambos %
métodos logramos obtener el total de alimentos generados por el campo de cultivo; $\mathbf{18}$ \textbf{Alimentos}

%%%%%%%%%%%%%%%%%%%%%%%%%%%%%%%%%%%%%%%
%%%%% Problema 1
%%%%%%%%%%%%%%%%%%%%%%%%%%%%%%%%%%%%%%

\begin{probgrupal}{}{}{}

\noindent Juan Carlos, siguiendo las órdenes de Patito ha comenzado a crear algunos campos de cultivo. Encuentre la cantidad total de alimentos que genera cada uno de ellos.

\begin{center}
\cultivos
\granja (g1) {$G_1$} ;
\granja (g2) at ($(g1) + (330:\grande)$) {$G_2$} ;
\granja (g3) at ($(g2) + (30:\grande)$) {$G_3$} ;
\finfin
\hspace{30mm}
\cultivos
\granja (g1) {$G_1$} ;
\granja (g2) at ($(g1) + (150:\grande)$) {$G_2$} ;
\granja (g3) [above of = g2] {$G_3$} ;
\granja (g4) [above of = g3] {$G_4$} ;
\granja (g5) at ($(g4) + (30:\grande)$) {$G_5$} ;
\granja (g6) at ($(g5) + (330:\grande)$) {$G_6$} ;
\granja (g7) [below of = g6] {$G_7$} ;
\granja (g8) at ($(g7) + (330:\grande)$) {$G_8$} ;
\finfin
\end{center}

\end{probgrupal}

\begin{solgrupal}
\noindent Agregar
\end{solgrupal}

%%%%%%%%%%%%%%%%%%%%%%%%%%%%%%%
%%%%% Problema 2
%%%%%%%%%%%%%%%%%%%%%%%%%%%%%%%

\begin{probgrupal}{}{}{}

Jonathan, siguiendo el ejemplo de Juan Carlos se dirigió a los alrededores de la futura colonia \emph{Colonia Champiñón Segunda}, construyendo los siguientes campos de cultivo. %
Encuentre la cantidad total de alimentos que genera cada uno de ellos.

\begin{center}
\cultivos
\granja (g1) {$G_1$} ;
\granja (g2) at ($(g1) + (150:\grande)$) {$G_2$} ;
\granja (g3) [above of = g2] {$G_3$} ;
\granja (g4) [above of = g3] {$G_4$} ;
\granja (g5) at ($(g4) + (30:\grande)$) {$G_5$} ;
\granja (g6) at ($(g5) + (330:\grande)$) {$G_6$} ;
\granja (g7) [below of = g6] {$G_7$} ;
\granja (g8) at ($(g7) + (330:\grande)$) {$G_8$} ;
\finfin
\hspace{15mm}
\cultivos
\granja (g1) {$G_1$} ;
\granja (g2) at ($(g1) + (30:\grande)$) {$G_2$} ;
\granja (g3) at ($(g1) + (330:\grande)$) {$G_3$} ;
\granja (g4) at ($(g3) + (30:\grande)$) {$G_4$} ;
\granja (g5) at ($(g3) + (330:\grande)$) {$G_5$} ;
\finfin
\hspace{15mm}
\cultivos
\granja (g1) {$G_1$} ;
\granja (g2) at ($(g1) + (30:\grande)$) {$G_2$} ;
\granja (g3) at ($(g1) + (330:\grande)$) {$G_3$} ;
\granja (g4) at ($(g3) + (30:\grande)$) {$G_4$} ;
\granja (g5) at ($(g3) + (330:\grande)$) {$G_5$} ;
\granja (g6) at ($(g5) + (30:\grande)$) {$G_6$} ;
\finfin
\end{center}

\end{probgrupal}


\begin{solgrupal}
Agregar
\end{solgrupal}

%%%%%%%%%%%%%%%%%%%%%%%%%%%
%%%%%%%%% Problema 3
%%%%%%%%%%%%%%%%%%%%%%%%%%%

\begin{probgrupal}{}{}{}
\noindent Siguiendo la misión de Patito, el equipo académico en su totalidad se dirigió a una zona llena de osbtáculos para crear granjas espaciales. 
Karla, una de sus integrantes, creó los siguientes tres campos de cultivo. En ellos, los hexágonos sombreados representan zonas en donde \textbf{NO hay %
una granja espacial}. Determine la cantidad de alimentos en total que genera cada campo de cultivo.

\begin{center}
\cultivos
\granja (g1) {$G_1$} ;
\granjaN (n2) at ($(g1) + (30:\grande)$) {} ;
\granjaN (n3) at ($(g1) + (330:\grande)$) {} ;
\granja (g2) at ($(n2) + (330:\grande)$) {$G_2$} ;
\granja (g3) at ($(g2) + (30:\grande)$) {$G_3$} ;
\granjaN (n4) [below of = g3] {} ;
\granja (g4) at ($(n4) + (30:\grande)$) {$G_4$} ;
\granja (g5) [above of = g4] {$G_5$} ;
\finfin
\hspace{20mm}
\cultivos
\granjaN (n1) {} ;
\granja (g1) [above of = n1] {$G_1$} ;
\granja (g2) [above of = g1] {$G_2$} ;
\granja (g3) [below of = n1] {$G_3$} ;
\granja (g4) [below of = g3] {$G_4$} ;
\granja (g5) at ($(n1) + (30:\grande)$) {$G_5$} ;
\granja (g6) at ($(g5) + (30:\grande)$) {$G_6$} ;
\granja (g7) at ($(n1) + (330:\grande)$) {$G_7$} ;
\granja (g8) at ($(g7) + (330:\grande)$) {$G_8$} ;
\granja (g9) at ($(n1) + (210:\grande)$) {$G_9$} ;
\granja (g10) at ($(g9) + (210:\grande)$) {$G_{10}$} ;
\granja (g11) at ($(n1) + (150:\grande)$) {$G_{11}$} ;
\granja (g12) at ($(g11) + (150:\grande)$) {$G_{12}$} ;
\finfin
\hspace{20mm}
\cultivos
\granja (g1) {$G_1$} ;
\granja (g2) at ($(g1) + (30:\grande)$) {$G_2$} ;
\granja (g3) at ($(g1) + (330:\grande)$) {$G_3$} ;
\granja (g4) at ($(g3) + (30:\grande)$) {$G_4$} ;
\granja (g5) at ($(g3) + (330:\grande)$) {$G_5$} ;
\finfin
\end{center}

\end{probgrupal}

\begin{solgrupal}
Agregar
\end{solgrupal}

%%%%%%%%%%%%%%%%%%%%%%%%%%%
%%%%%% Problema 4
%%%%%%%%%%%%%%%%%%%%%%%%%%

\begin{probgrupal}{}{}{}

Antoine, que trabajaba en la \emph{Colonia Champiñón Tercera}, ha diseñado los siguientes campos de cultivo. Determine la cantidad de alimento en total %
que entrega cada uno de ellos.

\begin{center}
\cultivos
\granja (g1) {$G_1$} ;
\granja (g2) at ($(g1) + (150:\grande)$) {$G_2$} ;
\granja (g3) [above of = g2] {$G_3$} ;
\granja (g4) [above of = g3] {$G_4$} ;
\granja (g5) at ($(g4) + (30:\grande)$) {$G_5$} ;
\granja (g6) at ($(g5) + (330:\grande)$) {$G_6$} ;
\granja (g7) [below of = g6] {$G_7$} ;
\granja (g8) at ($(g7) + (330:\grande)$) {$G_8$} ;
\finfin
\hspace{15mm}
\cultivos
\granja (g1) {$G_1$} ;
\granja (g2) at ($(g1) + (30:\grande)$) {$G_2$} ;
\granja (g3) at ($(g1) + (330:\grande)$) {$G_3$} ;
\granja (g4) at ($(g3) + (30:\grande)$) {$G_4$} ;
\granja (g5) at ($(g3) + (330:\grande)$) {$G_5$} ;
\granja (g6) at ($(g5) + (30:\grande)$) {$G_6$} ;
\finfin
\end{center}

\end{probgrupal}

\begin{solgrupal}
Agregar
\end{solgrupal}


%%%%%%%%%%%%%%%%%%%%%%
%%%%%% Problema 5
%%%%%%%%%%%%%%%%%%%%%%

\begin{probgrupal}{}{}{}

Jorge estaba explorando el planeta junto a su mascota capibara \emph{Capix}. Tras una extensa jornada, dio con una zona ideal para %
el desarrollo de granjas espaciales. Determine la cantidad total de alimentos generados por cada campo de cultivo respectivamente.

\begin{center}
\cultivos
\granja (g1) {$G_1$} ;
\granjaN (n2) at ($(g1) + (30:\grande)$) {} ;
\granjaN (n3) at ($(g1) + (330:\grande)$) {} ;
\granja (g2) at ($(n2) + (330:\grande)$) {$G_2$} ;
\granja (g3) at ($(g2) + (30:\grande)$) {$G_3$} ;
\granjaN (n4) [below of = g3] {} ;
\granja (g4) at ($(n4) + (30:\grande)$) {$G_4$} ;
\granja (g5) [above of = g4] {$G_5$} ;
\finfin
\hspace{20mm}
\cultivos
\granja (g1) {$G_1$} ;
\granjaN (n1) [above of = g1] {} ;
\granja (g11) [above of = n1] {$G_{11}$} ;
\granjaN (n2) [above of = g11] {} ;
\granja (g2) [above of = n2] {$G_2$} ;
\granja (g4) at ($(g1) + (30:\grande)$) {$G_4$} ;
\granja (g3) at ($(g1) + (150:\grande)$) {$G_3$} ;
\granja (g5) [above of = g4] {$G_5$} ;
\granja (g10) [above of = g5] {$G_{10}$} ;
\granja (g7) [above of = g10] {$G_7$} ;
\granja (g6) [above of = g3] {$G_6$} ;
\granja (g9) [above of = g6] {$G_9$} ;
\granja (g8) [above of = g9] {$G_8$} ;
\finfin
\end{center}

\end{probgrupal}

\begin{solgrupal}
Agregar
\end{solgrupal}

%%%%%%%%%%%%%%%%%%%%
%%%%%%%% Problema 6
%%%%%%%%%%%%%%%%%%%

\begin{probgrupal}{}{}{}

Catherina y Karla se encontraban en el vehículo de exploración \emph{Rover Cuack} deplazándose en los alrededores de la \emph{Colonia Pan de Ajo}. %
En su búsqueda, encontraron dos sitios aptos para la cosntrucción de granjas espaciales. Tras hacer ambos diseños, Patito les informa que solo uno de %
ellos podrá ser construido. Determine cuál de los siguientes campos de cultivo genera una \textbf{mayor cantidad} de alimentos. \emph{Los hexágonos sombreados %
NO tienen granjas}

\begin{center}
\cultivos
\granja (g1) {$G_1$} ;
\granja (g2) at ($(g1) + (30:\grande)$) {$G_2$} ;
\granja (g3) at ($(g1) + (150:\grande)$) {$G_3$} ;
\granja (g4) [above of = g1] {$G_4$} ;
\granja (g5) [above of = g4] {$G_5$} ;
\finfin
\hspace{20mm}
\cultivos
\granja (g1) {$G_1$} ;
\granja (g2) [above of = g1] {$G_2$} ;
\granja (g3) at ($(g2) + (30:\grande)$) {$G_3$} ;
\granja (g4) at ($(g3) + (330:\grande)$) {$G_4$} ;
\granja (g5) [below of = g4] {$G_5$} ;
\granja (g6) at ($(g5) + (210:\grande)$) {$G_6$} ;
\granjaN (n1) [above of = g6] {} ;
\finfin
\end{center}
\end{probgrupal}

\begin{solgrupal}
Agregar
\end{solgrupal}

%%%%%%%%%%%%%%%%%%%%%%%
%%%%%%%  Problema 7
%%%%%%%%%%%%%%%%%%%%%%

\begin{probgrupal}{}{}{}

Hein ha recibido la tarea de diseñar un campo de cultivo para la colonia de \emph{Champiñón Sagrado Morado Tercero}. Tras una tarde de %
arduo trabajo, Hein ha entregado a Patito dos diseños diferentes. Determine cuál de los siguientes campos de cultivos genera una \textbf{mayor cantidad} %
de alimentos. \emph{Los hexágonos sombreados NO tienen granjas.}

\begin{center}
\cultivos
\granja (g1) {$G_1$} ;
\granja (g2) at ($(g1) + (30:\grande)$) {$G_2$} ;
\granja (g3) at ($(g1) + (330:\grande)$) {$G_3$} ;
\granjaN (n1) at ($(g3) + (30:\grande)$) {} ;
\granja (g4) [above of = n1] {$G_4$} ;
\granja (g5) [below of = n1] {$G_5$} ;
\granja (g6) at ($(g5) + (30:\grande)$) {$G_6$} ;
\granja (g7) [above of = g6] {$G_7$} ;
\granja (g8) at ($(g6) + (30:\grande)$) {$G_8$} ;
\finfin
\hspace{20mm}
\cultivos
\granjaN (n1) {} ;
\granja (g1) [above of = n1] {$G_1$} ;
\granja (g2) [above of = g1] {$G_2$} ;
\granja (g3) [below of = n1] {$G_3$} ;
\granja (g4) [below of = g3] {$G_4$} ;
\granja (g5) at ($(n1) + (30:\grande)$) {$G_5$} ;
\granja (g6) at ($(g5) + (30:\grande)$) {$G_6$} ;
\granja (g7) at ($(n1) + (330:\grande)$) {$G_7$} ;
\granja (g8) at ($(g7) + (330:\grande)$) {$G_8$} ;
\granja (g9) at ($(n1) + (210:\grande)$) {$G_9$} ;
\granja (g10) at ($(g9) + (210:\grande)$) {$G_{10}$} ;
\granja (g11) at ($(n1) + (150:\grande)$) {$G_{11}$} ;
\granja (g12) at ($(g11) + (150:\grande)$) {$G_{12}$} ;
\finfin
\end{center}


\end{probgrupal}


\begin{solgrupal}
Agregar
\end{solgrupal}


%%%%%%%%%%%%%%%%%%%%%%%
%%%%  Problema 8
%%%%%%%%%%%%%%%%%%%%%%%

\begin{probgrupal}{}{}{}

Nabil y Hernán se encontraban analizando mapas de algunas zonas recientmente exploradas del planeta \emph{Champiñon}. Al hacerlo, descubireron una zona %
de tierras muy fértiles, ideal para construir algunas granjas. Tras evaluar diferentes opciones, al llegar a su nave presentaron tres opciones a Patito. %
Determine cuál de los siguientes campos de cultivo genera una \textbf{mayor cantidad de alimentos}. \emph{Los hexágonos sombreados NO tienen granjas}.

\begin{center}
\cultivosP
\granja (g1) {$G_1$} ;
\granja (g2) at ($(g1) + (30:\enano)$) {$G_2$} ;
\granja (g3) at ($(g1) + (330:\enano)$) {$G_3$} ;
\granjaN (n1) at ($(g3) + (30:\enano)$) {} ;
\granja (g4) [above of = n1] {$G_4$} ;
\granja (g5) [below of = n1] {$G_5$} ;
\granja (g6) at ($(g5) + (30:\enano)$) {$G_6$} ;
\granja (g7) [above of = g6] {$G_7$} ;
\granja (g8) at ($(g6) + (30:\enano)$) {$G_8$} ;
\finfin
\hspace{20mm}
\cultivosP
\granjaN (n1) {} ;
\granja (g1) [above of = n1] {$G_1$} ;
\granja (g2) [above of = g1] {$G_2$} ;
\granja (g3) [below of = n1] {$G_3$} ;
\granja (g4) [below of = g3] {$G_4$} ;
\granja (g5) at ($(n1) + (30:\enano)$) {$G_5$} ;
\granja (g6) at ($(g5) + (30:\enano)$) {$G_6$} ;
\granja (g7) at ($(n1) + (330:\enano)$) {$G_7$} ;
\granja (g8) at ($(g7) + (330:\enano)$) {$G_8$} ;
\granja (g9) at ($(n1) + (210:\enano)$) {$G_9$} ;
\granja (g10) at ($(g9) + (210:\enano)$) {$G_{10}$} ;
\granja (g11) at ($(n1) + (150:\enano)$) {$G_{11}$} ;
\granja (g12) at ($(g11) + (150:\enano)$) {$G_{12}$} ;
\finfin
\hspace{20mm}
\cultivosP
\granja (g1) {$G_1$} ;
\granjaN (n1) [above of = g1] {} ;
\granjaN (n2) [above of = n1] {} ;
\granja (g2) [above of = n2] {$G_2$} ;
\granja (g3) at ($(g1) + (330:\enano)$) {$G_3$} ;
\granjaN (n3) [above of = g3] {} ;
\granja (g4) [above of = n3] {$G_4$} ;
\granjaN (n4) [above of = g4] {} ;
\granja (g5) [above of = n4] {$G_5$} ;
\granja (g6) at ($(g3) + (30:\enano)$) {$G_6$} ;
\granja (g7) [above of = g6] {$G_7$} ;
\granjaN (n5) [above of = g7] {} ;
\granja (g7) [above of = n5] {$G_8$} ;
\granja (g8) at ($(g6) + (30:\enano)$) {$G_9$} ;
\granja (g9) [above of = g8] {$G_{10}$} ;
\granja (g10) [above of = g9] {$G_{11}$} ;
\granja (g11) at ($(g8) + (30:\enano)$) {$G_{12}$} ;
\finfin
\end{center}


\end{probgrupal}


\begin{solgrupal}
Agregar
\end{solgrupal}

\vspace{20mm}

\noindent Finalmente, tras una larga espera llegó el día de recibir a las nuevas colonias de humanos en el planeta \emph{Champiñon}. Patito y su equipo estaban muy emocionados con esta noticia, %
además de ogullosos de su trabajo preparando granjas espaciales para generar suficientes alimentos. Al momento de observar las naves colonia llegar y aterrizar sobre la superficie del planeta %
\emph{Champiñon}, un nuevo inconveniente cruzó la mente de Patito. Si bien el equipo había trabajado con mucho esfuerzo en crear las granjas espaciales, así como las reglas para su funcionamiento, %
no habían pensado en las reglas que debían seguir las colonias recién llegadas para adaptarse a las granjas. De esta forma, mientras continuaban aterrizando las naves espaciales, Patito llamó al %
equipo académico de CMAT a una reunión de emergencia. En ella, crearon un nuevo sistema para regular y organizar el funcionamiento de las colonias.

\section*{Colonias Espaciales}

\noindent Designaremos a las colonias espaciales dentro del hexágono con una $\mathbf{C}$. Además, una característica de cada una de estas colonias, será el \emph{\textbf{tamaño}}. El \emph{tamaño} de una colonia espacial %
corresponde a la \textbf{cantidad mínima de comida} que requiere la colonia para sobrevivir. Representaremos este tamaño sobre la letra $\mathbf{C}$ como se aprecia en los ejemplos inferiores.

\noindent En ellos, podemos identificar una colonia de tamaño $10$. Luego dos colonias, de tamaños $10$ y $12$. Y finalmente tres colonias, de tamaños $8$, $14$, $6$ y $2$, respectivamente.

\begin{center}
\cultivos
\granja (c1) {$\overset{\mathbf{10}}{C}$};
\finfin
\hspace{20mm}
\cultivos
\granja (c1) {$\overset{\mathbf{10}}{C}_1$} ;
\granja (c2) at ($(c1) +  (30:\grande)$) {$\overset{\mathbf{12}}{C}_2$} ;
\finfin
\hspace{20mm}
\cultivos
\granja (c1) {$\overset{\mathbf{8}}{C}_1$} ;
\granja (c2) at ($(c1) + (30:\grande)$) {$\overset{\mathbf{14}}{C}_2$} ;
\granja (c3) at ($(c1) + (330:\grande)$) {$\overset{\mathbf{6}}{C}_3$} ;
\granja (c4) at ($(c3) + (30:\grande)$) {$\overset{\mathbf{2}}{C}_4$} ;
\finfin
\end{center}

\vspace{10mm}

\noindent De esta forma, diremos que una colonia se encuentra \emph{\textbf{abastecida}}, si recibe la cantidad \textbf{exacta} de alimento que pide su \emph{tamaño}. Para ello, Patito y su equipo definieron una serie de reglas %
que determinan cómo se distribuye el alimento generado por las granjas espaciales entre las colonias.

\vspace{10mm}

\subsection*{Distribución de alimentos; Reglas}

\begin{itemize}
\item Cada \emph{colonia espacial} puede recibir alimentos \textbf{única y exclusivamente} de sus granjas espaciales \emph{\textbf{adyacentes}}.
\item La cantidad \emph{\textbf{mínima}} de alimento que podrá recibir una \emph{colonia espacial}, es igual a su \emph{tamaño}
\end{itemize}